%% =====================================================================
%%  Rapport de fin de projet - Master 1 Informatique et Mobilite
%%  Projet CVGen
%%  Zakariae ABDOUNI & Rania MAMOU - Encadrant : M. Frederic Cordier
%%  UHA - 2025/2026
%%
%%  Compilation : pdflatex RAPPORT_FIN_PROJET.tex (deux fois pour la
%%  table des matieres). Le logo est optionnel (logo_uha.png/.pdf).
%% =====================================================================
\documentclass[12pt,a4paper]{report}

\usepackage[utf8]{inputenc}
\usepackage[T1]{fontenc}
\usepackage[french]{babel}
\usepackage{mathptmx}
\usepackage[top=2.5cm,bottom=2.5cm,left=2.5cm,right=2.5cm,headheight=15pt]{geometry}
\usepackage[table]{xcolor}
\usepackage{graphicx}
\usepackage{tabularx}
\usepackage{array}
\usepackage{booktabs}
\usepackage{titlesec}
\usepackage{fancyhdr}
\usepackage{enumitem}
\usepackage{tikz}
\usetikzlibrary{shapes.geometric, arrows.meta, positioning}

%% ─── Couleurs (avant hyperref) ───────────────────────────────────
\definecolor{navyblue}{RGB}{31,41,55}
\definecolor{accentblue}{RGB}{37,99,235}
\definecolor{lightblue}{RGB}{219,234,254}
\definecolor{rowalt}{RGB}{246,248,251}
\definecolor{subtlegray}{RGB}{107,114,128}

\usepackage{hyperref}
\hypersetup{colorlinks=true, linkcolor=navyblue, urlcolor=accentblue}

\setlength{\parindent}{1em}
\setlength{\parskip}{6pt}
\renewcommand{\arraystretch}{1.3}

\titleformat{\chapter}[display]
  {\bfseries\color{navyblue}}
  {\filright\large CHAPITRE \thechapter}{4pt}{\Huge}
  [{\color{accentblue}\titlerule[1.5pt]}]
\titlespacing*{\chapter}{0pt}{0pt}{20pt}
\titleformat{\section}{\large\bfseries\color{navyblue}}{\thesection}{0.7em}{}
\titleformat{\subsection}{\normalsize\bfseries\color{navyblue}}{\thesubsection}{0.6em}{}

\pagestyle{fancy}
\fancyhf{}
\renewcommand{\headrulewidth}{0.8pt}
\renewcommand{\headrule}{{\color{accentblue}\hrule width\headwidth height\headrulewidth}}
\fancyhead[L]{\small\textbf{Projet CVGen}~\textcolor{subtlegray}{| Rapport de fin de projet}}
\fancyfoot[C]{\small\textcolor{subtlegray}{\thepage}}

\newcolumntype{L}[1]{>{\raggedright\arraybackslash}p{#1}}
\newcommand{\THead}[1]{\textcolor{white}{\textbf{#1}}}
\newcommand{\rowhead}{\rowcolor{navyblue}}

%% Capture d'écran optionnelle : affiche l'image si elle existe, sinon un
%% cadre gris indiquant le fichier à fournir. À remplacer par vos captures.
%% Usage : \capture{largeur}{nom_fichier}{légende}{label}
\newcommand{\capture}[4]{%
  \begin{figure}[h]\centering
  \IfFileExists{#2}{\fbox{\includegraphics[width=#1]{#2}}}{%
    \fbox{\begin{minipage}[c][4.5cm][c]{#1}\centering
      \color{subtlegray}\small\itshape
      [ Capture d'\'ecran \`a ins\'erer ]\\[4pt]
      \texttt{#2}
    \end{minipage}}}%
  \caption{#3}\label{#4}
  \end{figure}}

\begin{document}

%% ═══════════════════════════════════════════════════════════════
%% PAGE DE GARDE
%% ═══════════════════════════════════════════════════════════════
\begin{titlepage}
\thispagestyle{empty}
\centering

\IfFileExists{logo_uha.png}{\includegraphics[width=6.5cm]{logo_uha.png}\par}{%
  \IfFileExists{logo_uha.pdf}{\includegraphics[width=6.5cm]{logo_uha.pdf}\par}{%
    {\large\bfseries\color{navyblue}Universit\'e de Haute-Alsace\par}
    {\color{subtlegray}Facult\'e des Sciences et Techniques\par}}}

\vspace{1.2cm}
{\large\color{subtlegray}Master 1 --- Informatique et Mobilit\'e\par}
\vspace{0.3cm}
{\color{subtlegray}Projet de fin d'ann\'ee\par}

\vspace{1.6cm}
{\color{accentblue}\rule{\linewidth}{2pt}}\par
\vspace{0.6cm}
{\fontsize{46}{52}\selectfont\bfseries\color{navyblue}CVGen\par}
\vspace{0.4cm}
{\Large\color{navyblue}Une plateforme pour g\'en\'erer des CV\\ et postuler plus facilement en ligne\par}
\vspace{0.6cm}
{\color{accentblue}\rule{\linewidth}{2pt}}\par

\vfill

\begin{tabular}{rl}
\textbf{R\'ealis\'e par :} & Zakariae ABDOUNI \\
                           & Rania MAMOU \\[6pt]
\textbf{Encadr\'e par :}   & M. Fr\'ed\'eric Cordier \\[6pt]
\textbf{Fili\`ere :}       & Informatique et Mobilit\'e \\[6pt]
\textbf{Ann\'ee :}         & 2025 -- 2026 \\
\end{tabular}

\vspace{1.5cm}
\end{titlepage}

%% ═══════════════════════════════════════════════════════════════
%% REMERCIEMENTS
%% ═══════════════════════════════════════════════════════════════
\chapter*{Remerciements}
\addcontentsline{toc}{chapter}{Remerciements}

Nous tenons d'abord \`a remercier notre encadrant, M. Fr\'ed\'eric Cordier, pour
ses conseils et son suivi tout au long de ce projet. Nous remercions \'egalement
l'ensemble de l'\'equipe p\'edagogique du Master Informatique et Mobilit\'e de
l'Universit\'e de Haute-Alsace. Enfin, ce projet a \'et\'e r\'ealis\'e en bin\^ome et
nous avons beaucoup appris l'un de l'autre en travaillant ensemble.

%% ═══════════════════════════════════════════════════════════════
%% RÉSUMÉ
%% ═══════════════════════════════════════════════════════════════
\chapter*{R\'esum\'e}
\addcontentsline{toc}{chapter}{R\'esum\'e}

Ce rapport pr\'esente le projet \textbf{CVGen}, que nous avons d\'evelopp\'e dans le
cadre de notre Master~1. L'id\'ee de d\'epart est partie d'un constat simple : quand
on cherche un emploi ou une alternance, on passe beaucoup de temps \`a refaire son
CV pour chaque offre et \`a ressaisir les m\^emes informations sur chaque site de
recrutement.

Pour r\'epondre \`a ce probl\`eme, nous avons construit une application compos\'ee de
trois parties : un \textbf{serveur} (backend) qui g\`ere les comptes, les profils et
qui fait appel \`a une intelligence artificielle pour g\'en\'erer un CV adapt\'e \`a une
offre ; un \textbf{site web} o\`u l'utilisateur remplit son profil et g\'en\`ere ses
documents ; et une \textbf{extension pour le navigateur Chrome} qui remplit
automatiquement les formulaires de candidature sur les sites des entreprises.

Dans la suite, nous expliquons le contexte et les objectifs, l'analyse que nous
avons faite, la conception de l'application, les technologies que nous avons
utilis\'ees, ce que nous avons r\'eellement r\'ealis\'e, et surtout les nombreuses
difficult\'es que nous avons rencontr\'ees ainsi que la fa\c{c}on dont nous les avons
r\'esolues.

\vspace{4pt}
\textbf{Mots-cl\'es :} CV, candidature, ATS, intelligence artificielle,
Spring Boot, React, extension Chrome.

\tableofcontents
\thispagestyle{fancy}

%% ═══════════════════════════════════════════════════════════════
%% CHAPITRE 1 — INTRODUCTION
%% ═══════════════════════════════════════════════════════════════
\chapter{Introduction}

\section{Contexte}
Aujourd'hui, la plupart des candidatures se font en ligne. Les grandes
entreprises utilisent des logiciels appel\'es ATS (\textit{Applicant Tracking
System}) pour recevoir et trier automatiquement les CV. Ces logiciels lisent le
CV, cherchent des mots-cl\'es et classent les candidats avant m\^eme qu'un recruteur
les regarde.

En tant qu'\'etudiants qui postulons r\'eguli\`erement (stages, alternances), nous
avons remarqu\'e deux choses g\^enantes. D'une part, un CV qui n'est pas bien
structur\'e ou qui ne contient pas les bons mots-cl\'es peut \^etre \'ecart\'e tr\`es t\^ot.
D'autre part, chaque site de recrutement (SuccessFactors, Workday, Taleo,
SmartRecruiters, etc.) demande de retaper encore et encore les m\^emes
informations : nom, exp\'eriences, formations, comp\'etences\ldots C'est long et
d\'ecourageant.

\section{Probl\'ematique}
La question \`a laquelle nous avons voulu r\'epondre est donc la suivante : comment
aider un candidat \`a obtenir un CV adapt\'e \`a une offre pr\'ecise, et lui \'eviter de
ressaisir tout son profil sur chaque site, sans aller \`a l'encontre des r\`egles de
s\'ecurit\'e des navigateurs ?

\section{Objectifs}
Nous nous sommes fix\'e les objectifs suivants :
\begin{itemize}[leftmargin=1.4em, itemsep=2pt]
  \item permettre \`a l'utilisateur de cr\'eer un profil complet (exp\'eriences,
  formations, comp\'etences, langues, certifications, projets) ;
  \item lui permettre d'importer un CV existant (PDF ou Word) pour gagner du
  temps ;
  \item g\'en\'erer un CV adapt\'e \`a une offre donn\'ee, ainsi qu'une lettre de
  motivation ;
  \item donner un score de compatibilit\'e avec l'offre (score ATS) avec des
  conseils ;
  \item remplir automatiquement les formulaires de candidature gr\^ace \`a
  l'extension, et y d\'eposer le CV et la lettre.
\end{itemize}

\section{Plan du rapport}
Apr\`es cette introduction, le chapitre~2 pr\'esente l'application dans son
ensemble. Le chapitre~3 d\'ecrit les besoins. Le chapitre~4 explique la conception
et l'architecture. Le chapitre~5 liste les technologies. Le chapitre~6 d\'etaille
ce que nous avons r\'ealis\'e. Le chapitre~7 revient sur les difficult\'es
rencontr\'ees. Le chapitre~8 parle des tests et de l'organisation, et le
chapitre~9 conclut avec le bilan et les pistes d'am\'elioration.

%% ═══════════════════════════════════════════════════════════════
%% CHAPITRE 2 — PRÉSENTATION
%% ═══════════════════════════════════════════════════════════════
\chapter{Pr\'esentation du projet}

\section{Vue d'ensemble}
CVGen est compos\'e de trois parties qui communiquent entre elles autour d'une
m\^eme API :
\begin{itemize}[leftmargin=1.4em, itemsep=3pt]
  \item \textbf{le backend} : c'est la partie centrale. Il g\`ere les comptes, le
  profil, les appels \`a l'IA et la cr\'eation des documents ;
  \item \textbf{le site web} : c'est l\`a que l'utilisateur cr\'ee son compte,
  remplit son profil, g\'en\`ere ses CV et regarde son score ATS ;
  \item \textbf{l'extension Chrome} : elle travaille directement dans le
  navigateur, sur la page de candidature, pour remplir le formulaire \`a partir du
  profil.
\end{itemize}

\section{Qui utilise l'application}
L'utilisateur principal est le candidat. En pratique, il cr\'ee un compte, remplit
son profil (ou l'importe depuis un ancien CV), g\'en\`ere un CV pour une offre,
regarde son score, t\'el\'echarge le PDF, puis va sur le site de l'entreprise et
laisse l'extension remplir le formulaire \`a sa place.

\capture{0.85\linewidth}{captures/accueil.png}{Page d'accueil du site CVGen.}{fig:accueil}

\section{Ce que fait l'application}
Le tableau~\ref{tab:perimetre} r\'esume les fonctionnalit\'es que nous avons mises
en place.

\begin{table}[h]
\centering\small
\begin{tabularx}{\linewidth}{L{4.2cm}X}
\toprule
\rowhead \THead{Partie} & \THead{Ce qu'elle permet} \\
\midrule
\rowcolor{rowalt} Compte & Inscription, connexion (avec jetons de session). \\
Profil & Ajouter / modifier / supprimer ses exp\'eriences, formations, comp\'etences, langues, certifications et projets. \\
\rowcolor{rowalt} Import de CV & Lire un CV PDF ou Word et remplir le profil automatiquement. \\
G\'en\'eration & Cr\'eer un CV adapt\'e \`a une offre et une lettre de motivation, puis exporter en PDF. \\
\rowcolor{rowalt} Score ATS & Donner une note sur 100, les mots-cl\'es trouv\'es ou manquants, et des conseils. \\
Extension & Remplir les formulaires des sites de recrutement et y d\'eposer les documents. \\
\bottomrule
\end{tabularx}
\caption{R\'esum\'e des fonctionnalit\'es.}
\label{tab:perimetre}
\end{table}

%% ═══════════════════════════════════════════════════════════════
%% CHAPITRE 3 — BESOINS
%% ═══════════════════════════════════════════════════════════════
\chapter{Analyse des besoins}

Avant de coder, nous avons list\'e ce dont l'application avait besoin, en
s\'eparant les besoins fonctionnels (ce qu'elle doit faire) des besoins non
fonctionnels (comment elle doit le faire).

\section{Besoins fonctionnels}
\begin{itemize}[leftmargin=1.4em, itemsep=2pt]
  \item cr\'eer un compte et se connecter de fa\c{c}on s\'ecuris\'ee ;
  \item g\'erer un profil professionnel complet ;
  \item importer un CV pour pr\'e-remplir le profil ;
  \item g\'en\'erer un CV adapt\'e \`a une offre ;
  \item g\'en\'erer une lettre de motivation ;
  \item calculer un score ATS avec des conseils ;
  \item exporter les documents en PDF ;
  \item remplir automatiquement les formulaires et y joindre les documents.
\end{itemize}

\section{Besoins non fonctionnels}
\begin{itemize}[leftmargin=1.4em, itemsep=2pt]
  \item \textbf{S\'ecurit\'e} : les mots de passe ne doivent jamais \^etre stock\'es en
  clair, et l'API doit \^etre prot\'eg\'ee.
  \item \textbf{Rapidit\'e} : la g\'en\'eration et le remplissage doivent rester
  rapides pour ne pas \^etre p\'enibles.
  \item \textbf{Robustesse} : l'IA renvoie parfois des r\'esultats bizarres, il
  faut que l'application ne plante pas pour autant.
  \item \textbf{Compatibilit\'e} : les sites de recrutement sont tous diff\'erents,
  l'extension doit s'adapter \`a beaucoup de cas.
  \item \textbf{Code clair} : nous voulions un code organis\'e pour pouvoir le
  reprendre et le faire \'evoluer.
\end{itemize}

%% ═══════════════════════════════════════════════════════════════
%% CHAPITRE 4 — CONCEPTION
%% ═══════════════════════════════════════════════════════════════
\chapter{Conception et architecture}

\section{Architecture g\'en\'erale}
L'application suit une organisation classique client--serveur. Le site web et
l'extension sont les clients, et ils discutent tous les deux avec le m\^eme
serveur (l'API REST). Le serveur, lui, parle \`a la base de donn\'ees et \`a l'IA. La
figure~\ref{fig:archi} montre comment tout cela s'articule.

\begin{figure}[h]
\centering
\begin{tikzpicture}[
  node distance=0.9cm and 1.4cm,
  box/.style={rectangle, rounded corners, draw=navyblue, fill=lightblue,
              text width=3.2cm, align=center, minimum height=1.2cm, font=\small},
  ext/.style={rectangle, rounded corners, draw=accentblue, fill=white,
              text width=3.2cm, align=center, minimum height=1.2cm, font=\small},
  db/.style={cylinder, shape border rotate=90, aspect=0.3, draw=navyblue,
             fill=rowalt, minimum height=1.4cm, minimum width=2.2cm, align=center, font=\small},
  arr/.style={-{Stealth[length=2mm]}, thick, draw=subtlegray}
]
  \node[box] (web) {Site web\\ \scriptsize React};
  \node[ext, below=of web] (chrome) {Extension Chrome};
  \node[box, right=2.2cm of web, yshift=-0.9cm, fill=navyblue, text=white] (api)
       {Backend\\ \scriptsize API REST + JWT};
  \node[db, right=of api] (pg) {PostgreSQL};
  \node[ext, above=of api, fill=lightblue, draw=navyblue] (gem)
       {IA Gemini};
  \node[ext, below=1.2cm of api] (ats) {Sites ATS\\ \scriptsize (recruteurs)};

  \draw[arr] (web) -- (api);
  \draw[arr] (chrome) -- (api);
  \draw[arr] (api) -- (pg);
  \draw[arr] (api) -- (gem);
  \draw[arr] (chrome) -- (ats) node[midway, left, font=\scriptsize] {remplissage};
\end{tikzpicture}
\caption{Sch\'ema g\'en\'eral de l'application.}
\label{fig:archi}
\end{figure}

\section{Le backend}
Nous avons d\'ecoup\'e le backend par \og domaines \fg{}, c'est-\`a-dire par grandes
fonctionnalit\'es : l'authentification, le profil, la g\'en\'eration et le score ATS.
Chaque domaine est lui-m\^eme s\'epar\'e en trois couches : les contr\^oleurs (qui
re\c{c}oivent les requ\^etes), les services (qui contiennent la logique), et la partie
technique (base de donn\'ees, s\'ecurit\'e, appel \`a l'IA). Nous avons choisi cette
organisation pour que le code reste lisible et que chaque partie ait une
responsabilit\'e claire.

Un point important : nous ne renvoyons jamais directement les objets de la base
de donn\'ees au client. Nous passons toujours par des objets interm\'ediaires (DTO)
et une r\'eponse au format uniforme, ce qui \'evite d'exposer des informations
sensibles.

\section{La base de donn\'ees}
Les donn\'ees sont stock\'ees dans PostgreSQL, et nous utilisons des fichiers de
migration (Flyway) pour cr\'eer et faire \'evoluer les tables. Le
tableau~\ref{tab:entites} pr\'esente les principales tables.

\begin{table}[h]
\centering\small
\begin{tabularx}{\linewidth}{L{3.6cm}X}
\toprule
\rowhead \THead{Table} & \THead{Contenu} \\
\midrule
\rowcolor{rowalt} Utilisateurs & Comptes (email, mot de passe chiffr\'e, r\^ole). \\
Profils & Le profil li\'e \`a chaque utilisateur. \\
\rowcolor{rowalt} Exp\'eriences / Formations & Le parcours professionnel et scolaire. \\
Comp\'etences / Langues & Comp\'etences (avec niveau) et langues. \\
\rowcolor{rowalt} Certifications / Projets & Certifications et projets personnels. \\
CV g\'en\'er\'es & Les CV cr\'e\'es par l'IA. \\
\bottomrule
\end{tabularx}
\caption{Principales tables de la base.}
\label{tab:entites}
\end{table}

\section{L'intelligence artificielle}
C'est la partie qui nous a le plus int\'eress\'es. Nous utilisons le mod\`ele de
Google, \textbf{Gemini}, \`a travers son API. Pour ne pas r\'ep\'eter le m\^eme code
partout, nous avons cr\'e\'e un seul \og point d'entr\'ee \fg{} vers l'IA, qui g\`ere les
appels, les erreurs et les limites d'usage. Ensuite, quatre fonctionnalit\'es
l'utilisent :
\begin{itemize}[leftmargin=1.4em, itemsep=2pt]
  \item la \textbf{g\'en\'eration de CV} : on envoie \`a l'IA l'offre et le profil, et
  elle choisit puis reformule les \'el\'ements les plus pertinents ;
  \item la \textbf{lettre de motivation} : l'IA r\'edige une lettre personnalis\'ee ;
  \item le \textbf{score ATS} : l'IA compare le CV \`a l'offre et donne une note avec
  des mots-cl\'es et des conseils ;
  \item l'\textbf{import de CV} : l'IA transforme le texte brut d'un ancien CV en
  donn\'ees rang\'ees dans le profil.
\end{itemize}

Un probl\`eme que nous avons vite rencontr\'e : l'IA ne renvoie pas toujours un
texte propre (elle ajoute parfois des virgules en trop ou des balises). Nous
avons donc \'ecrit un petit outil pour nettoyer ses r\'eponses avant de les lire.

\section{La s\'ecurit\'e}
Pour la connexion, nous utilisons des jetons JWT. \`A la connexion, l'utilisateur
re\c{c}oit un jeton de courte dur\'ee et un jeton de rafra\^ichissement plus long. \`A
chaque requ\^ete, le serveur v\'erifie le jeton. Les mots de passe sont chiffr\'es
avec BCrypt et seuls le site web et l'extension sont autoris\'es \`a appeler l'API.

\section{L'extension Chrome}
L'extension respecte la derni\`ere version du format des extensions
(Manifest~V3). Elle est compos\'ee d'un programme de fond (\textit{service
worker}), de scripts inject\'es dans les pages des sites de recrutement, et d'une
petite fen\^etre (popup) pour la piloter. Les scripts inject\'es sont s\'epar\'es selon
leur r\^ole : d\'etecter les champs, faire la correspondance avec le profil, remplir,
g\'erer les sections r\'ep\'etables, les dates, l'extraction de l'offre et l'envoi des
fichiers.

%% ═══════════════════════════════════════════════════════════════
%% CHAPITRE 5 — TECHNOLOGIES
%% ═══════════════════════════════════════════════════════════════
\chapter{Technologies utilis\'ees}

Le tableau~\ref{tab:techno} regroupe les principales technologies que nous avons
utilis\'ees. Nous avons essay\'e de choisir des outils r\'epandus, bien document\'es,
que nous pouvions apprendre dans le temps imparti.

\begin{table}[h]
\centering\small
\begin{tabularx}{\linewidth}{L{3.4cm}X}
\toprule
\rowhead \THead{Cat\'egorie} & \THead{Technologies} \\
\midrule
\rowcolor{rowalt} Backend & Java 17, Spring Boot, Maven. \\
Base de donn\'ees & PostgreSQL, Flyway. \\
\rowcolor{rowalt} IA & Google Gemini (via son API). \\
Documents & Apache PDFBox et Apache POI (lecture des CV), g\'en\'eration de PDF. \\
\rowcolor{rowalt} S\'ecurit\'e & JWT, BCrypt, Spring Security. \\
Site web & React, TypeScript, Vite, Tailwind CSS. \\
\rowcolor{rowalt} Extension & Chrome Manifest V3, JavaScript. \\
Outils & Docker, Git, GitHub Actions. \\
\bottomrule
\end{tabularx}
\caption{Technologies du projet.}
\label{tab:techno}
\end{table}

%% ═══════════════════════════════════════════════════════════════
%% CHAPITRE 6 — RÉALISATION
%% ═══════════════════════════════════════════════════════════════
\chapter{R\'ealisation}

\section{Le backend}
Nous avons d\'evelopp\'e l'API morceau par morceau. D'abord l'authentification
(inscription, connexion), puis la gestion du profil avec toutes ses
sous-parties, ensuite l'import de CV, et enfin la g\'en\'eration et le score ATS.
Nous avons aussi ajout\'e une page de documentation automatique de l'API pour
nous y retrouver, et un endroit unique qui g\`ere les erreurs afin de renvoyer des
messages clairs (par exemple quand l'IA a atteint sa limite d'utilisation).

\section{Le site web}
Le site permet de s'inscrire, de se connecter, de remplir son profil (avec une
petite barre qui indique le pourcentage de compl\'etion), d'importer un CV via une
fen\^etre d\'edi\'ee, et enfin de g\'en\'erer un CV pour une offre puis de voir son score
ATS et de t\'el\'echarger le PDF. Nous avons soign\'e l'interface pour qu'elle reste
simple \`a utiliser.

\capture{0.85\linewidth}{captures/profil.png}{\'Edition du profil sur le site web.}{fig:profil}

\capture{0.85\linewidth}{captures/score_ats.png}{Affichage du score ATS apr\`es g\'en\'eration d'un CV.}{fig:ats}

\section{L'extension Chrome}
C'est la partie la plus difficile techniquement, et celle sur laquelle nous avons
pass\'e le plus de temps. L'extension sait :
\begin{itemize}[leftmargin=1.4em, itemsep=2pt]
  \item rep\'erer et remplir les champs sur de nombreux sites de recrutement, en
  reconnaissant les champs gr\^ace \`a leurs libell\'es ;
  \item g\'erer les sections que l'on peut ajouter plusieurs fois (plusieurs
  exp\'eriences, plusieurs formations\ldots) en cliquant toute seule sur les
  boutons \og Ajouter \fg{} ;
  \item s'adapter aux composants compliqu\'es de certains sites (par exemple les
  listes d\'eroulantes de SuccessFactors) ;
  \item d\'eposer automatiquement le CV et la lettre dans les champs de fichier.
\end{itemize}

\capture{0.55\linewidth}{captures/extension_popup.png}{La fen\^etre (popup) de l'extension Chrome.}{fig:popup}

\capture{0.85\linewidth}{captures/formulaire_rempli.png}{Exemple de formulaire de candidature rempli automatiquement par l'extension.}{fig:rempli}

%% ═══════════════════════════════════════════════════════════════
%% CHAPITRE 7 — DIFFICULTÉS
%% ═══════════════════════════════════════════════════════════════
\chapter{Difficult\'es rencontr\'ees}

Cette partie est sans doute la plus int\'eressante, car c'est l\`a que nous avons le
plus appris. Nous d\'ecrivons les principaux probl\`emes et comment nous nous en
sommes sortis.

\section{Les \'ev\'enements \og non humains \fg{}}
Pour des raisons de s\'ecurit\'e, les navigateurs marquent tout ce qui est fait par
un script comme \og non humain \fg{}. Certains sites (comme SuccessFactors) refusent
alors le remplissage de leurs listes d\'eroulantes. Nous avons d\^u tester
plusieurs m\'ethodes pour contourner cela, et quand vraiment aucune ne marchait,
nous avons choisi de surligner en orange les champs \`a finir \`a la main plut\^ot que
de laisser l'utilisateur dans le flou.

\section{Les champs cach\'es (Shadow DOM ferm\'e)}
Sur certains sites (SmartRecruiters), les champs sont enferm\'es dans une zone
inaccessible aux scripts normaux. Apr\`es pas mal de recherches, nous avons fini
par utiliser un outil plus puissant fourni par Chrome pour \og atteindre \fg{} ces
champs et y \'ecrire caract\`ere par caract\`ere, comme le ferait une vraie personne.

\section{Les champs qui se remplissent puis disparaissent}
C'est le probl\`eme qui nous a le plus \'enerv\'es. Sur certains sites, on voyait le
champ se remplir, puis se vider tout seul, ou l'entr\'ee dispara\^itre au moment de
sauvegarder. Nous avons compris qu'il fallait remplir en plusieurs fois, bien
envoyer les \'ev\'enements, et parfois cocher automatiquement une case (par exemple
\og poste en cours \fg{}) pour que le site accepte une exp\'erience sans date de fin.

\section{Les r\'eponses de l'IA}
L'IA ne donne pas toujours un r\'esultat parfait : parfois le texte n'est pas un
JSON valide, parfois le CV n'est pas assez adapt\'e \`a l'offre. Nous avons am\'elior\'e
nos consignes (le \og prompt \fg{}), nettoy\'e les r\'eponses, et fait en sorte que le
remplissage du formulaire continue m\^eme si la g\'en\'eration \'echoue, pour ne pas
tout bloquer.

\section{Les limites de l'API gratuite}
Comme nous utilisons la version gratuite de l'IA, nous \'etions souvent bloqu\'es
par des limites d'utilisation. Nous avons ajout\'e des nouvelles tentatives
automatiques et des messages d'erreur clairs pour comprendre ce qui se passait.

\section{Le vocabulaire diff\'erent selon les sites}
Un m\^eme niveau ne s'\'ecrit pas pareil partout : un niveau de langue \og Bon \fg{}
dans notre profil correspond \`a \og Avanc\'e \fg{} sur un site, et \og Bac+5 \fg{}
correspond ailleurs \`a \og Bac~+4/+5 \fg{}. Nous avons cr\'e\'e une petite table de
correspondances pour traduire automatiquement nos valeurs vers celles attendues
par chaque site.

%% ═══════════════════════════════════════════════════════════════
%% CHAPITRE 8 — TESTS ET ORGANISATION
%% ═══════════════════════════════════════════════════════════════
\chapter{Tests et organisation du travail}

\section{Comment nous avons test\'e}
Nous avons surtout test\'e \og \`a la main \fg{}, en essayant l'application sur de vrais
formulaires de candidature. Pour comprendre les bugs de l'extension, nous avons
ajout\'e beaucoup de messages dans la console (pr\'efix\'es par \texttt{[CVGen]}), ce
qui nous a permis de voir exactement o\`u \c{c}a coin\c{c}ait \`a chaque essai. C'est
gr\^ace \`a ces messages que nous avons pu corriger la plupart des probl\`emes du
chapitre pr\'ec\'edent. Nous avons aussi mis en place une int\'egration continue
(GitHub Actions) qui v\'erifie que le projet compile bien \`a chaque modification.

\section{R\'epartition du travail}
Nous avons travaill\'e en bin\^ome, en avan\c{c}ant par \og sprints \fg{} d'une semaine.
La r\'epartition s'est faite naturellement :
\begin{itemize}[leftmargin=1.4em, itemsep=2pt]
  \item \textbf{Zakariae} s'est occup\'e surtout du backend (base de donn\'ees,
  s\'ecurit\'e, IA, g\'en\'eration, score ATS) et du site web ;
  \item \textbf{Rania} s'est occup\'ee de l'extension Chrome (remplissage des
  formulaires, sections dynamiques, envoi des documents) et de la liaison avec
  l'API.
\end{itemize}
Nous avons utilis\'e Git pour travailler chacun de notre c\^ot\'e sans nous g\^ener, et
nous \'ecrivions un petit rapport chaque semaine pour suivre l'avancement.

%% ═══════════════════════════════════════════════════════════════
%% CHAPITRE 9 — BILAN
%% ═══════════════════════════════════════════════════════════════
\chapter{Bilan et perspectives}

\section{Bilan}
Au final, nous sommes plut\^ot contents du r\'esultat. L'application fait ce que
nous voulions : on peut cr\'eer son profil, importer un CV, g\'en\'erer un CV adapt\'e \`a
une offre, voir un score ATS et laisser l'extension remplir les formulaires. Tout
n'est pas parfait, mais l'essentiel fonctionne et l'extension marche sur une
douzaine de sites de recrutement.

\section{Ce que nous avons appris}
Ce projet nous a fait progresser sur beaucoup de points :
\begin{itemize}[leftmargin=1.4em, itemsep=2pt]
  \item organiser un projet et s\'eparer le code proprement ;
  \item utiliser une IA dans une application (et g\'erer ses d\'efauts) ;
  \item d\'evelopper \`a la fois c\^ot\'e serveur (Java/Spring) et c\^ot\'e site
  (React) ;
  \item programmer une extension de navigateur, ce que nous n'avions jamais
  fait ;
  \item travailler \`a deux avec Git sans \'ecraser le travail de l'autre.
\end{itemize}

\section{Ce qu'on pourrait am\'eliorer}
Il reste des choses \`a faire si nous avions plus de temps :
\begin{itemize}[leftmargin=1.4em, itemsep=2pt]
  \item ajouter de vrais tests automatiques (nous avons surtout test\'e \`a la
  main) ;
  \item cr\'eer un tableau de bord pour suivre les candidatures envoy\'ees ;
  \item proposer plusieurs mod\`eles de CV ;
  \item supporter encore plus de sites de recrutement ;
  \item g\'erer d'autres langues que le fran\c{c}ais.
\end{itemize}

\section{Conclusion}
Ce projet nous a permis de partir d'un vrai besoin (postuler plus facilement) et
d'aller jusqu'\`a une application qui fonctionne, en passant par toutes les \'etapes :
r\'efl\'echir au probl\`eme, concevoir, coder, tester et corriger. Nous avons surtout
appris en nous confrontant \`a des probl\`emes concrets, parfois fastidieux, mais
toujours formateurs. C'est une exp\'erience que nous garderons pour la suite de nos
\'etudes et de notre vie professionnelle.

\end{document}
